\section{Strategy as a Game}
Game theory models strategic interactions between firms as games, where each
firm's payoff depends on its own actions and those of its competitors. Key
concepts include:

\begin{itemize}
    \item \textbf{Nash Equilibrium:} A stable state where no player can improve their payoff by
          unilaterally changing their strategy.
    \item \textbf{Dominant Strategy:} A strategy that yields a higher payoff regardless of the
          opponent's actions.
    \item \textbf{Mixed Strategy:} Randomizing over multiple strategies to keep opponents
          uncertain.
\end{itemize}

\textbf{Example:} In a price war, if both firms lower prices, they may end up in a Nash
Equilibrium with low profits. If one firm can commit to a high price (eg.,
through a loyalty program), it can gain a competitive advantage.

\subsection{Formalising Games}
\textbf{Actors:} Firms or players in the game. \\
\textbf{Actions:} Possible strategies or moves each firm can make. \\
\textbf{Payoffs:} Outcomes or profits resulting from the combination of actions taken
by all firms. \\
\textbf{Timing:} Simultaneous (e.g., price setting) vs. sequential (e.g., entry
decisions).

Here is the completed section for your cheatsheet, incorporating the game
theory models and case analysis from your course materials.

\subsection{Finding Optimal Strategies}
\textbf{Backward Induction:} A technique for solving sequential games by working
backward from the final move to the initial decision to determine the optimal
sequence of actions.

\textbf{Participation Constraint:} A player will only enter or remain in a partnership
if their expected payoff exceeds their opportunity cost.

\textbf{Added Value:} The value a firm brings to a partnership, calculated as:
$Value_{WithFirm} - Value_{WithoutFirm}$.

\subsection{Market Entry and Deterrence}
Firms can use strategic actions to prevent competitors from entering a market:
\begin{itemize}
    \item \textbf{Limit Pricing:} Setting a low price (e.g., \$10) to make entry unprofitable for
          competitors facing high fixed costs (e.g., \$1300 factory setup). \item \textbf{Overproduction:} Expanding capacity beyond current demand to signal an
          aggressive response to entry. \item \textbf{Niche Differentiation:} Carving out specific segments (e.g., Mattel's focus on
          sports games) to avoid direct head-to-head rivalry with a dominant incumbent.
\end{itemize}

\subsection{Co-opetition: Value Creation \& Capture}
\textbf{Co-opetition:} A strategy where firms simultaneously cooperate to build total
value and compete to capture a larger share of that value.

\textbf{Example: Apple vs. Google (Early Smartphone Era)}
\begin{itemize}
    \item \textbf{Cooperation:} Google provided essential apps (Maps, Search) to increase the
          iPhone's total value ($Value \approx \$200$).
    \item \textbf{Competition:} As Google launched Android, the relationship shifted. Google's
          opportunity cost rose (partnerships with Samsung), allowing it to capture more
          value from Apple.
\end{itemize}

\section{Case Study: The Golden Age of Video Games}
\subsection{Atari: The Rise and Fall (1972--1985)}
\begin{itemize}
    \item \textbf{Success Factors:} First-mover advantage with simple, intuitive games (Pong).
          Used a ``Razor-and-Blade'' model: low margins on consoles, high margins on
          cartridges.
    \item \textbf{The Failure (Atari Shock 1983):} A flood of poor-quality, uninspired software
          from independent developers (e.g., Activision) caused a market collapse.
    \item \textbf{Strategic Error:} Refusal to cannibalize its own products (``eat its own
          babies'') allowed competitors to innovate faster.
\end{itemize}

\subsection{Nintendo: The Monopoly Strategy (1985--1991)}
Nintendo revived the industry by implementing a ``closed'' ecosystem.
\begin{itemize}
    \item \textbf{Strategic Control:} Use of a \textbf{Security Chip} (lock-out) ensured only Nintendo-approved cartridges could run
          on the NES.
    \item \textbf{Network Effects:} Nintendo leveraged \textbf{Two-Sided Platforms}—attracting more users to bring in more high-quality
          developers (Super Mario, Zelda), which in turn attracted more users.
    \item \textbf{Strict Licensing:} Restricted licensees to 5 titles/year and mandated exclusive
          2-year clauses to prevent competitors (Atari, Coleco) from gaining software.
    \item \textbf{Contrived Shortage:} Carefully managed inventory to keep demand and prices
          high.
\end{itemize}

\textbf{Strategic Takeaway:} While Nintendo achieved a sustained competitive advantage
through organizational control (VRIO), its ``iron-grip'' led to significant
legal and partner resentment.
